\documentclass[11pt]{article}
\usepackage[a4paper,margin=27mm]{geometry}
\usepackage{amsmath,amssymb,amsthm,mathtools}
\usepackage[T1]{fontenc}
\usepackage[utf8]{inputenc}
\usepackage{lmodern}
\usepackage{microtype}
\usepackage{xcolor}
\usepackage{hyperref}
\hypersetup{colorlinks=true,linkcolor=blue!55!black,citecolor=blue!55!black,urlcolor=blue!55!black}
\setlength{\parindent}{0pt}
\setlength{\parskip}{5pt plus 1pt}
\newtheorem{theorem}{Theorem}[section]
\newtheorem{lemma}[theorem]{Lemma}
\newtheorem{proposition}[theorem]{Proposition}
\theoremstyle{remark}
\newtheorem{remark}[theorem]{Remark}
\newcommand{\per}{\operatorname{per}}
\newcommand{\one}{\mathbf 1}
\newcommand{\RR}{\mathbb R}
\newcommand{\QQ}{\mathbb Q}

\title{Dittert's conjecture in dimension 9\\
\large An exact computer-assisted boundary-exclusion proof}
\author{Pedro Paulo Marques do Nascimento}
\date{23 July 2026}

\begin{document}
\maketitle

\begin{center}
\fbox{\parbox{0.90\textwidth}{\textbf{Status.} This is an exact, reproducible
working proof. It has not yet undergone independent peer review. The finite
dimension-specific calculations are checked by two separately implemented
standard-library verifiers.}}

\small Licensed under \href{https://creativecommons.org/licenses/by/4.0/}{CC BY 4.0}.
\end{center}

\begin{abstract}
For a nonnegative $n\times n$ matrix $A$ whose entries sum to $n$, let $r_i,c_j$
be its row and column sums and put
\[
 \Phi(A)=\prod_i r_i+\prod_j c_j-\per A.
\]
We prove Dittert's conjecture in dimension $n=9$.  A maximal Li scaling
reduces the boundary problem to whole, marginal, and proper tight cuts.
The proper cuts are excluded by exact staircase-face bounds.  The direct
scalar estimate for a marginal cut has a small failure interval.  On that
interval, stationarity forces a large column containing the minimum row;
cofactor nonnegativity bounds every other row meeting this column.  Seven
exact degree-82 Bernstein certificates then exclude every possible column
support.  The proof uses only a two-independent-zero permanent bound and does
not assume a three-zero permanent-minimum theorem.
\end{abstract}

\section{Statement and published inputs}

Let
\[
 K_n=\left\{A=(a_{ij})\in\RR_{\ge0}^{n\times n}:
       \sum_{i,j}a_{ij}=n\right\},
 \qquad
 \gamma_n=\frac{n!}{n^n}.
\]
For $A\in K_n$, write
\[
 R=\prod_{i=1}^n r_i,\qquad
 C=\prod_{j=1}^n c_j,\qquad
 P=\per A.
\]
Thus $\Phi(A)=R+C-P$ and $\Phi(J_n/n)=2-\gamma_n$.

\begin{theorem}\label{thm:main}
For every $A\in K_9$,
\[
 \Phi(A)\le 2-\frac{9!}{9^9},
\]
with equality if and only if $A=J_9/9$.
\end{theorem}

We use the following published results.
\begin{enumerate}
\item Hwang proved that the only possible positive $\Phi$-maximizer in $K_n$
is $J_n/n$ \cite{Hwang1986}.
\item Cheon and Wanless proved that a partly decomposable matrix cannot
maximize $\Phi$, and that a maximizer cannot have its complete zero set equal,
up to row and column permutations, to one proper rectangular block
\cite[Theorems 3.2 and 3.7]{CheonWanless2012}.
\item Li's cut characterization says that a nonnegative $n\times n$ matrix
$S$ is doubly superstochastic if and only if
\[
 \sum_{i\in I,j\in J}s_{ij}\ge |I|+|J|-n
\]
for all $I,J\subseteq[n]$ \cite{Li1990}; see also
\cite[Lemma 2.2]{CheonWanless2012}.
\item On a staircase face of the Birkhoff polytope, the barycenter minimizes
the permanent \cite{Hwang1985}.
\item Minc's averaging theorem for permanent minimizers with repeated support
rows and columns gives the block form used below \cite[Theorem 1]{Minc1984};
see also \cite[Section 2]{PulaSongWanless2011}.
\item The van der Waerden permanent theorem gives $\per B\ge\gamma_n$ for
every doubly stochastic $B$, with equality only at $J_n/n$
\cite{Egorychev1981,Falikman1981}.
\end{enumerate}
The support-stationarity equation and the large-column support argument used
below are proved directly; they are not additional external assumptions.

\section{Boundary structure and the shared deficit}

Because $K_n$ is compact, $\Phi$ has a global maximizer.  By Hwang's theorem,
a nonuniform maximizer must lie on the boundary.

\begin{lemma}[Two independent zeroes]\label{lem:independent-zeros}
Every boundary $\Phi$-maximizer has two zero entries in distinct rows and
distinct columns.
\end{lemma}
\begin{proof}
Suppose a nonempty zero set contains no such pair.  Fix a zero at $(i,j)$.
Every other zero shares row $i$ or column $j$.  If there were both a zero
$(i,j')$, $j'\ne j$, and a zero $(i',j)$, $i'\ne i$, those two zeroes would be
independent.  Hence all zeroes lie in one row or all lie in one column.

If that row or column is entirely zero, the matrix is partly decomposable.
Otherwise its complete zero set is one proper rectangular block.  Both cases
are excluded for a maximizer by Cheon--Wanless.
\end{proof}

Assume for contradiction that a boundary maximizer $A$ satisfies
$\Phi(A)\ge2-\gamma_n$.  Since the row sums and column sums each have
arithmetic mean $1$, AM--GM gives $R,C\le1$.  Consequently $P\le\gamma_n$.
Write
\[
 P=\gamma_n-\delta,\qquad 0\le\delta\le\gamma_n,
\]
and put $R=1-\rho$, $C=1-\sigma$.  The assumed lower bound gives
\begin{equation}\label{eq:shared-deficit}
 \rho+\sigma\le\delta.
\end{equation}
In particular,
\begin{equation}\label{eq:product-lower}
 R\ge1-\delta,\qquad C\ge1-\delta,
 \qquad RC\ge1-\delta.
\end{equation}

\section{The maximal Li scaling}

For $I,J\subseteq[n]$, let $s(A[I,J])$ be the sum of the indicated submatrix,
and define
\begin{equation}\label{eq:lambda}
 \lambda=
 \min_{\substack{I,J\subseteq[n]\\k:=|I|+|J|-n>0}}
 \frac{s(A[I,J])}{k}.
\end{equation}
A maximizer is fully indecomposable by the Cheon--Wanless exclusion.  Hence
every numerator in \eqref{eq:lambda} is positive, so $\lambda>0$; taking
$I=J=[n]$ shows $\lambda\le1$.

Li's theorem says that $A/\lambda$ is doubly superstochastic.  Thus there is a
doubly stochastic matrix $B$ such that
\begin{equation}\label{eq:domination}
 \lambda B\le A\quad\hbox{entrywise},
 \qquad P\ge\lambda^n\per B.
\end{equation}
Choose $I,J$ attaining \eqref{eq:lambda}, and set
\[
 u=n-|I|,\qquad v=n-|J|,\qquad k=n-u-v>0.
\]
There are three cases: the whole cut $(u,v)=(0,0)$, a marginal cut where
exactly one of $u,v$ is zero, and a proper cut where $u,v\ge1$.

If $u=v=0$, then $\lambda=1$.  Since $B\le A$ and both matrices have total
sum $n$, $A=B$ is doubly stochastic.  The van der Waerden theorem and
$P\le\gamma_n$ force $A=J_n/n$, contrary to the boundary assumption.

\section{The two-independent-zero permanent gap}

Let $\Omega_n^{(2)}$ be the face of the Birkhoff polytope specified by
$b_{12}=b_{21}=0$, and put $m=n-2$.  Minc's averaging theorem permits a
permanent minimizer on this face to be chosen in the form
\begin{equation}\label{eq:Y}
 Y=
 \begin{pmatrix}
 x_1&0&a_1\one^{\mathsf T}\\
 0&x_2&a_2\one^{\mathsf T}\\
 a_1\one&a_2\one&xJ_m
 \end{pmatrix},
 \qquad x_i=1-ma_i,
 \qquad a_1+a_2+mx=1.
\end{equation}
Here $0\le a_i\le1/m$.

\begin{proposition}[One-variable reduction]\label{prop:face-reduction}
For $m\ge4$, the minimum permanent in \eqref{eq:Y} occurs with $a_1=a_2$.
Consequently the minimum of $\per B$ on $\Omega_n^{(2)}$ is the minimum, on
\[
 \frac{m-2}{m^2}\le x\le\frac1m,
\]
of
\begin{equation}\label{eq:Pn}
 P_n(x)=m!x^{m-2}
 \left(x^2a^2+2mxab^2+m(m-1)b^4\right),
\end{equation}
where
\begin{equation}\label{eq:ab}
 a=\frac{2-m+m^2x}{2},\qquad b=\frac{1-mx}{2}.
\end{equation}
\end{proposition}
\begin{proof}
Put $s=a_1+a_2$, $p=a_1a_2$, and $x=(1-s)/m$.  Counting permanent terms in
\eqref{eq:Y} gives, apart from the positive factor $m!x^{m-2}$,
\[
 E_m(s,p)=x^2(1-ms+m^2p)
 +mx\bigl(s^2-(2+ms)p\bigr)+m(m-1)p^2.
\]
For fixed $s$,
\[
 \frac{\partial E_m}{\partial p}
 =2m(m-1)p+(m+1)s^2-ms-1.
\]
Since $p\le s^2/4$ and $0\le s\le2/m$, this is at most
\[
 D_m(s)=\frac{(m^2+m+2)s^2-2ms-2}{2}.
\]
The quadratic is convex, while
\[
 D_m(0)=-1,
 \qquad D_m(2/m)=-1+\frac2m+\frac4{m^2}<0
 \quad(m\ge4).
\]
Thus the derivative is strictly negative.  For fixed $s$, the permanent is
minimized by maximizing $p$, namely by $a_1=a_2=b$.  Solving the stochastic
constraints yields \eqref{eq:ab} and the stated interval.  A second count
gives \eqref{eq:Pn}.
\end{proof}

For $n=9$, set
\begin{equation}\label{eq:L}
 L_9=\frac{47533}{50000000}.
\end{equation}

\begin{proposition}[Exact two-zero bound]\label{prop:two-zero-bound}
Every $B\in\Omega_9^{(2)}$ satisfies
\[
 \per B>L_9>\gamma_9.
\]
\end{proposition}
\begin{proof}
Here $m=7$ and the complete interval is $[5/49,1/7]$.  The primary verifier
expands $P_9(x)-L_9$ over $\QQ$, verifies positive endpoint values, and uses
an exact rational Sturm sequence to prove that it has no root on the interval.

The independent verifier instead evaluates the matrix \eqref{eq:Y} with
\eqref{eq:ab} by the literal exact Ryser formula at ten rational points,
interpolates the degree-nine polynomial, and proves positivity by exact
Bernstein coefficients on 256 subintervals.  Finally,
\begin{equation}\label{eq:L-gap}
 L_9-\gamma_9=
 \frac{3348865477}{239148450000000}>0.
\end{equation}
\end{proof}

\section{Marginal tight cuts}\label{sec:marginal}

We first record the stationarity equation used twice below.

\begin{lemma}[Support stationarity]\label{lem:stationarity}
At a global maximizer, for every positive entry $a_{ij}$,
\begin{equation}\label{eq:stationarity}
 \frac{R}{r_i}+\frac{C}{c_j}-\per A(i\mid j)=\Phi(A),
\end{equation}
where $A(i\mid j)$ is obtained by deleting row $i$ and column $j$.
\end{lemma}
\begin{proof}
The left side is $\partial\Phi/\partial a_{ij}$.  Moving a small amount of
mass between positive entries shows that all partial derivatives on the
support have a common value $\mu$.  Euler's identity, using homogeneity of
degree $n$, gives
\[
 n\Phi(A)=\sum_{i,j}a_{ij}\frac{\partial\Phi}{\partial a_{ij}}
 =\mu\sum_{i,j}a_{ij}=n\mu.
\]
Thus $\mu=\Phi(A)$.
\end{proof}

Suppose the tight cut is marginal.  Its ratio in \eqref{eq:lambda} is an
average of a nonempty set of row sums or column sums.  By transposition, it is
enough to treat a minimum row.  Choose $i$ with
\[
 r_i=\lambda=1-a,\qquad 0\le a<1.
\]
Multiply \eqref{eq:stationarity} by $a_{ij}$ and sum over $j$.  Expansion of
the permanent along row $i$ gives
\[
 R+C\sum_j\frac{a_{ij}}{c_j}-P=r_i\Phi(A).
\]
Therefore
\begin{equation}\label{eq:qi}
 q_i:=\sum_j\frac{a_{ij}}{c_j}
 =1-(1+\theta)a,
 \qquad \theta=\frac{R-P}{C}.
\end{equation}
By \eqref{eq:product-lower} and $P=\gamma_n-\delta$,
\begin{equation}\label{eq:theta}
 \theta\ge R-P\ge1-\gamma_n.
\end{equation}
Let $c_{\max}=\max_jc_j$.  Since $q_i>0$, equations
\eqref{eq:qi}--\eqref{eq:theta} imply
\[
 \frac{1-a}{c_{\max}}\le q_i\le1-(2-\gamma_n)a.
\]
Thus the last denominator is positive and
\begin{equation}\label{eq:h}
 c_{\max}\ge1+h(a),
 \qquad h(a)=\frac{(1-\gamma_n)a}{1-(2-\gamma_n)a}.
\end{equation}

AM--GM on the other eight row sums and on the other eight column sums gives
\begin{align}
 R&\le f(a):=(1-a)\left(1+\frac a8\right)^8,
 \label{eq:f}\\
 C&\le g(h(a)):=(1+h(a))\left(1-\frac{h(a)}8\right)^8.
 \label{eq:g}
\end{align}
Define
\begin{equation}\label{eq:D}
 D(a)=2-f(a)-g(h(a)).
\end{equation}
The shared deficit implies $\delta\ge D(a)$.  Moreover, $f$ decreases for
$a>0$, $h$ increases, and $g$ decreases for positive arguments, so $D$ is
strictly increasing.  The exact check
\begin{equation}\label{eq:D-endpoint}
 D(3/100)>\gamma_9
\end{equation}
therefore confines every feasible value to $0\le a<3/100$.

The dominated doubly stochastic matrix $B$ inherits the two independent
zeroes of $A$, so Proposition~\ref{prop:two-zero-bound} gives
\begin{equation}\label{eq:marginal-permanent}
 P\ge L_9(1-a)^9.
\end{equation}
The elementary scalar contradiction is
\begin{equation}\label{eq:H}
 H(a):=L_9(1-a)^9-\gamma_9+D(a)>0.
\end{equation}
After clearing the positive denominator $8^8[1-(2-\gamma_9)a]^9$, its
numerator has degree 18.  The primary verifier proves by exact Sturm
sequences that $H>0$ on
\begin{equation}\label{eq:safe-intervals}
 [0,1/500]\ \cup\ [1/200,3/100].
\end{equation}
The independent verifier reconstructs this numerator from exact values and
proves the same two claims by Bernstein positivity.  In view of
\eqref{eq:L-gap}, only the interval
\begin{equation}\label{eq:risk-interval}
 \frac1{500}\le a\le\frac1{200}
\end{equation}
requires a stronger argument.

\subsection{The large-column support bound}

The weighted-average identity
\[
 \frac{q_i}{r_i}=\sum_{j:a_{ij}>0}\frac{a_{ij}}{r_i}\frac1{c_j}
\]
shows that some $j_*$ with $a_{ij_*}>0$ satisfies
\begin{equation}\label{eq:distinguished-column}
 c_{j_*}\ge\frac{r_i}{q_i}\ge1+h(a).
\end{equation}
The important extra fact, beyond \eqref{eq:h}, is that this large column
contains the minimum row.

Put
\begin{equation}\label{eq:q-support}
 d(a)=\left(1-\frac{h(a)}8\right)^8,
 \qquad q(a)=2-\gamma_9-d(a).
\end{equation}
For every $\ell$ with $a_{\ell j_*}>0$, cofactor nonnegativity and
Lemma~\ref{lem:stationarity} give
\[
 \frac{R}{r_\ell}
 =\Phi(A)-\frac{C}{c_{j_*}}+\per A(\ell\mid j_*)
 \ge2-\gamma_9-d(a)=q(a).
\]
Indeed, $C/c_{j_*}$ is the product of the other eight column sums, whose sum
is at most $8-h(a)$, so AM--GM bounds it by $d(a)$.  On
\eqref{eq:risk-interval}, $q$ is increasing and the verifier checks exactly
that $q(1/500)>1$.  Since $R\le1$, every row meeting $j_*$ satisfies
\begin{equation}\label{eq:row-cap}
 r_\ell\le\frac{R}{q(a)}\le\frac1{q(a)}<1.
\end{equation}

Let $z$ be the number of zero entries in column $j_*$.  Full
indecomposability excludes a column with only one positive entry, while a
full-support column is impossible because \eqref{eq:row-cap} would make the
sum of all row sums less than 9.  Hence $1\le z\le7$.

The support has $9-z$ rows: the distinguished row of sum $1-a$, and
$8-z$ other rows bounded by $1/q(a)$.  The remaining $z$ row sums have total
chosen to make the grand total 9.  AM--GM, followed by increasing each capped
row to its cap, gives
\begin{equation}\label{eq:Uz}
 R\le U_z(a):=(1-a)q(a)^{-(8-z)}
 \left(\frac{8+a-(8-z)/q(a)}{z}\right)^z.
\end{equation}
For completeness, the monotonicity in a capped variable $x$ follows by
differentiating the logarithm of
$x[(T-x)/z]^z$, where the uncapped total has average greater than $1$ while
$x<1$.

On the other hand, $\Phi(A)\ge2-\gamma_9$, \eqref{eq:g}, and
\eqref{eq:marginal-permanent} imply
\begin{equation}\label{eq:R0}
 R\ge R_0(a):=2-\gamma_9+L_9(1-a)^9-g(h(a)).
\end{equation}
It remains to check $R_0(a)>U_z(a)$ for $1\le z\le7$.

Here is the exact polynomial certificate.  Put
\begin{align*}
 e&=1-(2-\gamma_9)a,&
 t&=e-\frac{(1-\gamma_9)a}{8},\\
 Q&=(2-\gamma_9)e^8-t^8=e^8q(a),&
 S&=(2-\gamma_9+L_9(1-a)^9)e^9
     -(e+(1-\gamma_9)a)t^8=e^9R_0(a),\\
 W_z&=(8+a)Q-(8-z)e^8.
\end{align*}
All denominators are positive on \eqref{eq:risk-interval}, and
$R_0>U_z$ is equivalent to
\begin{equation}\label{eq:Gz}
 G_z(a):=z^zQ^8S
 -(1-a)e^{9+8(8-z)}W_z^z>0.
\end{equation}
Each $G_z$ has degree 82.  The primary verifier expands it directly and
converts it to the degree-82 Bernstein basis on $[1/500,1/200]$; every one of
the $7\cdot83=581$ coefficients is strictly positive.  The smallest is
approximately $3.7489581\cdot10^{-5}$ (the exact rational is retained by the
program and logged by the audit).  This contradicts
\eqref{eq:Uz}--\eqref{eq:R0} and excludes every marginal cut.

\section{Proper tight cuts}

Suppose now that $u,v\ge1$.  Define
\[
 \varepsilon_r=\sum_{i\in I^c}r_i-u,
 \qquad
 \varepsilon_c=\sum_{j\in J^c}c_j-v.
\]
For positive numbers $x_1,\dots,x_n$ of sum $n$ and product $X$, binary
Pinsker after grouping a subset and its complement gives
\[
 \left(\sum_{i\in S}x_i-|S|\right)^2
 \le-\frac n2\log X.
\]
Applying this to the row and column sums, then using Cauchy--Schwarz and
\eqref{eq:product-lower}, yields
\begin{equation}\label{eq:T}
 |\varepsilon_r|+|\varepsilon_c|
 \le\sqrt{-n\log(RC)}
 \le\sqrt{\frac{n\delta}{1-\delta}}=:T.
\end{equation}
The cut identity
\[
 s(A[I,J])=k-\varepsilon_r-\varepsilon_c+s(A[I^c,J^c])
\]
therefore implies
\begin{equation}\label{eq:lambda-proper}
 \lambda\ge1-\frac Tk.
\end{equation}
For $n=9$, the exact inequality $n\gamma_n<1-\gamma_n$ ensures $T<1\le k$.

Because $B$ is doubly stochastic,
\[
 s(B[I,J])=k+s(B[I^c,J^c])\ge k.
\]
Tightness and \eqref{eq:domination} give the reverse bound after
multiplication by $\lambda$, so equality holds and $B[I^c,J^c]=0$.  Thus $B$
lies on the staircase face with a prescribed $u\times v$ zero block.
Hwang's theorem gives
\begin{equation}\label{eq:nu}
 \per B\ge\nu_{n;u,v}
 :=\frac{\gamma_{n-u}\gamma_{n-v}}{\gamma_k},
 \qquad k=n-u-v,
\end{equation}
where $\gamma_0=1$.  The barycenter has
$(n-v)!(n-u)!/k!$ supported permutations, all of the same weight; both
verifiers reconstruct this count exactly.

Using \eqref{eq:lambda-proper}, Bernoulli's inequality, and
$1-\delta\ge1-\gamma_n$ gives
\begin{align*}
 \nu_{n;u,v}\left(1-\frac Tk\right)^n-(\gamma_n-\delta)
 &\ge \nu_{n;u,v}-\gamma_n+\delta
 -\frac{n\nu_{n;u,v}}k
  \sqrt{\frac{n\delta}{1-\gamma_n}}.
\end{align*}
Completing the square in $\sqrt\delta$ gives the uniform lower bound
\begin{equation}\label{eq:eta}
 \eta_{n;u,v}
 :=\nu_{n;u,v}-\gamma_n
 -\frac{n^3\nu_{n;u,v}^2}{4k^2(1-\gamma_n)}.
\end{equation}
Both verifiers check $\eta_{9;u,v}>0$ for all
\[
 u,v\ge1,\qquad u+v\le8,
\]
namely all 28 proper cuts.  The smallest value occurs at
$(u,v,k)=(1,7,1)$ (and symmetrically at $(7,1,1)$), where
\[
 \nu_{9;1,7}=\frac{315}{262144},
 \qquad
 \eta_{9;1,7}=
 \frac{9881810825836802335}{6282434825977483864571904}>0.
\]
This contradicts $P=\gamma_n-\delta$ in every proper-cut case.

The whole, marginal, and proper cases exhaust every tight cut.  No boundary
global maximizer exists.  Hwang's positive-support theorem identifies the
sole remaining maximizer as $J_9/9$, proving Theorem~\ref{thm:main}.

\section{Exact computational audit}

No floating-point value is used in a correctness decision.  The primary
standard-library verifier checks:
\begin{enumerate}
\item the bivariate two-zero equalization identity and its sign;
\item the degree-nine permanent polynomial and the lower bound \eqref{eq:L}
by a rational Sturm sequence;
\item denominator and endpoint signs, \eqref{eq:D-endpoint}, and zero Sturm
roots for the degree-18 numerator on both safe intervals;
\item $q(1/500)>1$ and all 581 exact Bernstein coefficients in
\eqref{eq:Gz};
\item the staircase barycenter count and all 28 values in \eqref{eq:eta}.
\end{enumerate}

The independent verifier imports no code from the primary verifier.  It uses
a literal exact Ryser sum for the two-zero matrices and Bernstein positivity
on 256 subintervals.  It reconstructs the marginal numerator from 19 exact
values and each support polynomial from 83 exact scalar values, checks
off-grid values, and proves the required positivity by independently
implemented Bernstein conversion.  Thus the principal expansions and both
Sturm computations are not shared.

The integration test runs both programs in ordinary and optimized Python mode
and requires identical output.  It also inflates $L_9$ to a false bound and
weakens the accounting of capped support rows; both independently mutated
programs must reject both changes.

\section{Scope and status}

This note establishes dimension 9.  Together with the other exact working
proofs in the accompanying repository, the covered dimensions are
\[
 4,5,9,10,11,12,13,14,15,16.
\]
Dimensions $6,7,8$ are not settled by these materials.  The present result
should be described as an exact computer-assisted working proof until the
mathematical reduction and both programs have undergone independent peer
review.

\section*{Reproducibility}

From the \texttt{n9} directory, run
\begin{verbatim}
python3 -I verify_dittert_n9.py
python3 -I audit_dittert_n9_stdlib.py
python3 -I test_n9_package.py
\end{verbatim}
The verifier final lines must be
\begin{verbatim}
CERTIFIED
INDEPENDENT AUDIT CERTIFIED
\end{verbatim}
and the tests must report \texttt{OK}.

\begin{thebibliography}{99}

\bibitem{CheonWanless2012}
G.-S. Cheon and I. M. Wanless,
\emph{Some results towards the Dittert conjecture on permanents},
Linear Algebra Appl. 436 (2012), 791--801.
\href{https://doi.org/10.1016/j.laa.2010.08.041}{doi:10.1016/j.laa.2010.08.041}.

\bibitem{Egorychev1981}
G. P. Egorychev,
\emph{Solution of the van der Waerden problem for permanents},
Soviet Math. Dokl. 23 (1981), 619--622.

\bibitem{Falikman1981}
D. I. Falikman,
\emph{A proof of van der Waerden's conjecture on the permanent of a stochastic matrix},
Mat. Zametki 29 (1981), 931--938.

\bibitem{Hwang1985}
S.-G. Hwang,
\emph{Minimum permanent on faces of staircase type of the polytope of doubly stochastic matrices},
Linear Multilinear Algebra 18 (1985), 271--306.
\href{https://doi.org/10.1080/03081088508817694}{doi:10.1080/03081088508817694}.

\bibitem{Hwang1986}
S.-G. Hwang,
\emph{A note on a conjecture on permanents},
Linear Algebra Appl. 76 (1986), 31--44.
\href{https://doi.org/10.1016/0024-3795(86)90212-0}{doi:10.1016/0024-3795(86)90212-0}.

\bibitem{Li1990}
C.-K. Li,
\emph{On certain convex matrix sets},
Discrete Math. 79 (1990), 323--326.

\bibitem{Minc1984}
H. Minc,
\emph{Minimum permanents of doubly stochastic matrices with prescribed zero entries},
Linear Multilinear Algebra 15 (1984), 225--243.
\href{https://doi.org/10.1080/03081088408817592}{doi:10.1080/03081088408817592}.

\bibitem{PulaSongWanless2011}
K. Pula, S.-Z. Song, and I. M. Wanless,
\emph{Minimum permanents on two faces of the polytope of doubly stochastic matrices},
Linear Algebra Appl. 434 (2011), 232--238.
\href{https://doi.org/10.1016/j.laa.2010.08.022}{doi:10.1016/j.laa.2010.08.022}.

\end{thebibliography}

\end{document}
